% P16.18: Hempel-Oppenheim 1948 Parts I-III, with complete logical proofs.
\chapter{Deductive Explanation, Laws and Relative Emergence}
\label{sec:p16-deductive-nomological-explanation}

A deductive explanation specifies a result, the general principles
used to obtain it, and the conditions under which those principles
apply. Its logical derivation, the truth of its premises and the
evidence available for those premises are separate questions.
Hempel and Oppenheim's 1948 account makes these distinctions
explicit and then develops a more restrictive formal criterion
than its familiar laws-plus-conditions summary
\cite[parts~I--III, pp.~135--164]{HempelOppenheim1948Explanation}.

This chapter presents that model and its original qualifications.
It does not identify the model with every conception of explanation.
The paper's quantitative account of systematic power receives
its separate treatment in \cref{sec:p15-systematic-power}.

\section{The explanatory question and the status of the premises}
\label{sec:p16-dn-question}

\begin{hbdefinition}{def:p16-dn-schema}{Explanandum, explanans and the deductive pattern}
The \emph{explanandum} \(E\) is a sentence describing the
phenomenon to be explained. The \emph{explanans} consists
of general principles \(L_1,\ldots,L_r\) and, when needed,
antecedent or accompanying conditions \(C_1,\ldots,C_k\).
The elementary deductive pattern is
\[
             L_1,\ldots,L_r,C_1,\ldots,C_k\entails E .
\]
The result sentence describes specified features of an occurrence,
not every property of an event or object.
In the 1948 account a sound explanation requires:
\begin{enumerate}
\item deduction of \(E\) from the explanans;
\item general laws that are necessary to the derivation;
\item empirical content of the explanatory account;
\item truth of its premises.
\end{enumerate}
These are the source's preliminary requirements R1--R4.
An empty list of particular conditions is allowed when the
general principles suffice. The more precise criterion below
concerns a singular \(E\), rather than explanation of one
general law by another.
\index{explanation!deductive-nomological}\index{explanandum}\index{explanans}
\end{hbdefinition}

The source regards prediction and explanation as sharing this
deductive form: their difference concerns when the result
is known and what the argument is being used to do.
This claim does not say that a derivation selects a causal
direction, that a premise was actually measured before
the event, or that all statistical explanation is deductive.
The paper explicitly leaves its statistical and inductive
extensions for separate work.

Truth is an ideal correctness condition in this model.
Evidence can warrant confidence in a premise without
establishing its truth conclusively. Revising the evidence
can therefore change confidence in an explanation without
making the same fixed proposition change from true to false.
The source's historical Bode-law example illustrates this
distinction; it is not a rule for assigning numerical
probabilities to theories.

An event's uniqueness does not prevent describing it
as an instance of specified repeatable characteristics.
Likewise a dependence on history does not logically
exclude a law whose conditions include that history.
Whether a useful empirical generalization is available
remains a separate question. The source's cotton-market
and linguistic examples are explicitly incomplete
explanatory sketches, not validated causal studies.

The source also treats familiarity and empathy as
possible aids to discovery rather than criteria of
explanatory correctness. Neither appears among the
formal conditions below. The distinct task evidence
about analogical discovery is supplied in
\cref{sec:p19-analogy-and-problem-transfer}.

\section{The exact model language and its law classes}
\label{sec:p16-dn-language}

\begin{hbdefinition}{def:p16-dn-language}{The equality-free relational language}
Use a finite or countable first-order relational signature
with individual constants, variables and finite-arity
predicates. There are no nonconstant function symbols
and no identity predicate. Domains are nonempty.
Sentences are closed formulas. Defined predicates are
expanded before applying the source's syntactic criteria.
Assume at least one ground atom is available.
Truth constants are the usual Boolean abbreviations
using such an atom. Constants and predicates displayed
in an example are part of that example's signature.

Write \(\Gamma\entails A\) for classical semantic
consequence and \(A\equiv B\) for logical equivalence.
The soundness and countable completeness results
\cref{thm:fol-soundness,cor:fol-completeness-countable}
identify this consequence relation with the source's
classical derivability relation. Truth in the intended
interpretation \(\mathfrak A\), written
\(\mathfrak A\models A\), is a different assertion.

A \emph{singular} sentence contains no variables and
is therefore a Boolean combination of ground atoms.
A \emph{basic} sentence is a ground atom or its negation.
A \emph{generalized} sentence has a nonempty prefix
of quantifiers followed by a quantifier-free matrix.
It is \emph{universal} when the prefix contains only
universal quantifiers. It is \emph{purely generalized},
or \emph{purely universal}, when it additionally has
no individual constants.
A sentence is \emph{essentially generalized} when
it is not equivalent to any singular sentence.
An essentially universal sentence must also be in
universal form.
\end{hbdefinition}

The intended predicates are assumed to be qualitative,
without an essential reference to a particular named
object or location. This is an interpretive assumption,
not a property established merely by removing constants
from the printed formula. For example, hiding a reference
to a particular basket inside a primitive predicate
does not remove that reference from its meaning.
The source explicitly leaves a fully general recognition
criterion for qualitative primitives unresolved.

\begin{hblemma}{thm:p16-dn-ground-language}{proved=annex:p16.dn-ground-language}{Ground freedom and genuine generality}
Every assignment of truth values to the ground atoms
of this relational language is realized by a structure
in which distinct constant names denote distinct objects.
Consequently entailment between singular sentences
is exactly their Boolean entailment.

The sentences
\[
 \forall x\,P(x),\qquad
 \forall x(P(x)\to Q(x)),\qquad
 \forall x\,R(a,x),\qquad
 \exists x\,P(x)
\]
are essentially generalized when the displayed predicate
symbols are distinct primitives of the stated arities.
In contrast,
\[
 \forall x(P(x)\lor P(a))\equiv P(a),\qquad
 \forall x(P(x)\lor\neg P(x))\equiv\top .
\]
Thus the mere presence of quantifiers does not establish
essential generality.
\end{hblemma}

\begin{hbdefinition}{def:p16-dn-law-theory}{Fundamental and derivative laws and theories}
With the intended interpretation fixed, the 1948
definitions are:
\begin{enumerate}
\item A fundamental lawlike sentence is purely universal.
      It is a fundamental law when it is also true in
      \(\mathfrak A\).
\item A derivative law is essentially universal but
      not purely universal, and follows from a set
      of fundamental laws.
\item A fundamental theory is a purely generalized
      sentence true in \(\mathfrak A\).
\item A derivative theory is essentially generalized
      but not purely generalized, and follows from
      a set of fundamental theories.
\end{enumerate}
A law or theory belongs to the corresponding union
of its fundamental and derivative classes.
These uses of \emph{theory} name individual sentences,
possibly combining several principles. They differ
from the general use of theory for a set of sentences.
\index{law!fundamental and derivative}\index{lawlike sentence}
\end{hbdefinition}

\begin{hbtheorem}{thm:p16-dn-law-hierarchy}{proved=annex:p16.dn-law-hierarchy}{The hierarchy and its exact qualifications}
Every law in this model is a theory, and every theory
is true in \(\mathfrak A\). The fundamental and
derivative classes of either kind are disjoint.
Fundamental laws can be logical truths or have
vacuous antecedents.

If \(\forall y\forall xR(y,x)\) is true, then
\(\forall xR(a,x)\) is a derivative law.
If \(\exists xP(x)\) is true, it is a fundamental
theory but is not in universal form.
Two essentially generalized sentences can have
a conjunction equivalent to a singular sentence:
with \(A=P(a)\) and
\(U=\forall xQ(x)\),
\[
                   (A\lor U)\land(A\lor\neg U)\equiv A .
\]
This is why essential generality must be checked
for the combined theoretical component, rather
than inferred from its separate premises.
\end{hbtheorem}

A true universal sentence about a fixed finite
collection can be accidental. Conversely a finite
extension in the actual world does not by itself
disqualify a qualitative universal. The source
distinguishes actual finiteness from restriction
built into a predicate's meaning. Its basket,
refrigerator and hidden-name examples concern this
interpretive issue; the formal language assumption
does not turn it into an automatic syntactic test.

\section{The final potential-explanation criterion}
\label{sec:p16-dn-verification}

\begin{hbdefinition}{def:p16-dn-potential-explanation}{The 1948 verification-witness criterion}
For singular sentences \(C,E\), the ordered pair
\((T,C)\) is a \emph{potential explanans} for \(E\)
when all of the following hold:
\begin{enumerate}
\item \(T\) is essentially generalized.
\item \(T,C\entails E\).
\item There exists a set \(\Gamma\) of basic sentences
      such that \(\{T\}\cup\Gamma\) is satisfiable,
      \(\Gamma\entails C\), and \(\Gamma\not\entails E\).
\end{enumerate}
This is the source's final condition (7.8).
The basic sentences provide a possible verification
of the initial conditions which does not, by itself,
verify the result. They are not required to be the
actual observations made in \(\mathfrak A\).

The pair is an \emph{explanans} in the source's
truth-qualified sense when it is potential,
\(T\) is a theory as defined above, and
\(\mathfrak A\models C\). A singular \(E\) is
explainable by \(T\) when some such \(C\) exists.
\index{explanation!potential}\index{explanation!verification witness}
\end{hbdefinition}

\begin{hbtheorem}{thm:p16-dn-verification-properties}{proved=annex:p16.dn-verification-properties}{Nonvacuity, truth and equivalence}
Every potential explanans satisfies:
\[
 T\land C\text{ is satisfiable},\qquad
 C\not\entails E,\qquad
 E\text{ is neither valid nor contradictory}.
\]
The theoretical component is necessary to the
deduction, although individual particular
conditions can be dispensable.
Every truth-qualified explanans has
\(\mathfrak A\models E\).

Replacing \(C,E\) by equivalent singular
sentences, or \(T\) by an equivalent essentially
generalized sentence, preserves potential
explanation. The witness's compatibility model
and its countermodel to \(E\) cannot be the same model.
Requiring \(\{T,\neg E\}\cup\Gamma\)
to be satisfiable would contradict the very
deduction the criterion demands.
\end{hbtheorem}

\begin{hbtheorem}{thm:p16-dn-finite-witness}{proved=annex:p16.dn-finite-witness}{A finite certificate drawn from the initial-condition atoms}
In the witness condition it suffices to consider
finite, consistent sets \(\Gamma\) containing
only literals whose atoms occur in \(C\).
If \(C\) contains \(m\) distinct ground atoms,
at most \(3^m\) such literal sets need be
considered: an atom can be asserted, denied
or left unassigned.

For each candidate, \(\Gamma\entails C\) and
\(\Gamma\not\entails E\) are finite truth-table
questions over the atoms of \(C,E\).
Satisfiability of \(\{T\}\cup\Gamma\) is a
first-order question. This finite certificate
description does not supply a general
first-order decision procedure or a decision
test for the qualitative interpretation of
the primitives.
\end{hbtheorem}

\section{Self-explanation and an irrelevant-law construction}
\label{sec:p16-dn-self-explanation}

\begin{hbtheorem}{thm:p16-dn-common-content}{proved=annex:p16.dn-common-content}{The complete decomposition behind source Theorem 7.7}
Suppose \(C,E\) are singular and \(T,C\entails E\).
Set
\[
 E_1=C\lor E,\qquad E_2=\neg C\lor E,\qquad C_1=C.
\]
Then
\[
 E\equiv E_1\land E_2,\qquad
 C\equiv C_1\land E_1,\qquad T\entails E_2 .
\]
The strongest common consequence of \(C\) and
\(E\) is their disjunction \(C\lor E\).
If they have no common consequence other
than logical truths, then \(T\entails E\).
Thus forbidding all shared logical content
would exclude every case in which the
particular conditions are indispensable.
\end{hbtheorem}

This decomposition concerns logical content,
not a claim that a scientist actually reasoned
in the displayed order. The source permits
partial sharing of content and instead adds
the basic-verification witness to exclude
a different defect.

\begin{hbtheorem}{thm:p16-dn-smuggled-result}{proved=annex:p16.dn-smuggled-result}{Why deduction plus law dependence is insufficient}
Let
\[
 T=\forall xP(x),\qquad C=P(a)\to R(a,b),
                         \qquad E=R(a,b).
\]
Then \(T\) is essentially generalized,
\(T,C\entails E\), and \(C\not\entails E\).
There are interpretations making both
\(T\) and \(C\) true. Nevertheless this
pair fails the final verification-witness
criterion.

More generally, if \(T\entails A\) for a
singular \(A\), and \(\neg A\land\neg E\)
is satisfiable, then \(C=A\to E\) satisfies
\(T,C\entails E\) but \(C\not\entails E\).
If \(E\) is true in the intended interpretation,
that \(C\) is also true there. The stated
satisfiability condition is essential;
this construction is not a claim about
every lawlike sentence, including tautologies.
\end{hbtheorem}

\begin{hbtheorem}{thm:p16-dn-worked-certificates}{proved=annex:p16.dn-worked-certificates}{Four complete certificate calculations}
The following examples use independent primitive
predicates of the displayed arities.
\begin{enumerate}
\item For \(T=\forall x(P(x)\to Q(x))\),
      \(C=P(a)\), \(E=Q(a)\), the literal
      set \(\{P(a)\}\) is a witness.
\item For \(T=\forall xP(x)\), \(C=\top\),
      \(E=P(a)\), the empty literal set
      is a witness. No initial fact is needed.
\item For the first \(T\), let
      \(C=R(a,b)\land P(a)\land U(b)\) and
      \(E=Q(a)\land R(a,b)\).
      The three conjunct literals witness
      potential explanation despite the
      shared \(R(a,b)\).
\item For \(T=\forall xP(x)\),
      \(C=P(a)\lor Q(a)\), \(E=P(a)\),
      the set \(\{Q(a)\}\) is a possible
      witness. In an intended interpretation
      with all \(P\)'s true and all \(Q\)'s
      false, \(T,C,E\) are true but that
      witness is false. The criterion
      therefore need not identify the
      verification actually available
      in the intended situation.
\end{enumerate}
\end{hbtheorem}

\section{Relative emergence and missing bridge principles}
\label{sec:p16-dn-emergence}

\begin{hbdefinition}{def:p16-dn-relative-emergence}{An occurrence relative to theory, parts and attributes}
Fix an intended structure \(\mathfrak A\), a whole
named \(w\), a part relation \(\mathsf{Pt}\),
and a selected attribute vocabulary \(G\).
Let \(D_{\mathsf{Pt},G}(w)\) be the specified
description of those parts using \(G\), including
the positive and negative facts supplied by
that description. Let \(\mathcal T\) be the
available set of theoretical sentences.

For an actual occurrence
\(\mathfrak A\models W(w)\), call it emergent
relative to these choices when
\[
               \mathcal T\cup D_{\mathsf{Pt},G}(w)
                                      \not\entails W(w).
\]
This is the source's relational explication,
with the available description made explicit.
For finitely many named parts and a finite
relational \(G\), a complete selected-attribute
description can give one signed ground atom
for every \(G\)-predicate and tuple of those
parts. It does not automatically describe
attributes omitted from \(G\).

A characteristic-level statement applies this
test to each of its actual occurrences.
If there are no such occurrences, the universal
statement is vacuous and establishes no
existence claim.
\index{emergence!relative to theory and vocabulary}
\end{hbdefinition}

\begin{hbtheorem}{thm:p16-dn-emergence-monotonicity}{proved=annex:p16.dn-emergence-monotonicity}{Dependence on the available information}
With \(W,w\) fixed, adding theoretical or descriptive
premises cannot destroy an existing derivation of
\(W(w)\). It can remove relative emergence.
This observation concerns entailment; an inconsistent
enlargement entails every sentence and supplies
no sound factual explanation.

If at least one part \(p\) is named, allowing as
a selected part attribute the whole-dependent
formula \(\mathsf{Pt}(p,w)\land W(w)\) trivializes
the test whenever that attribute is supplied
as true. Restricting the part vocabulary is
therefore substantive. The argument requires
an actual part and does not apply to a partless
case merely by vacuity.
\end{hbtheorem}

\begin{hbtheorem}{thm:p16-dn-emergence-examples}{proved=annex:p16.dn-emergence-examples}{Missing structure and missing theory are separate}
Use named objects \(p,q,w\), a unary property \(A\),
a binary connection predicate \(R\), and a
whole-property predicate \(W\). Suppose the
part membership is fixed by
\(\mathsf{Pt}(p,w),\mathsf{Pt}(q,w)\), the
selected attribute data initially give only
\(A(p),A(q)\), and the theory is
\[
             \forall x\forall y\forall z
 ((\mathsf{Pt}(x,z)\land\mathsf{Pt}(y,z)\land
                       A(x)\land A(y)\land R(x,y))\to W(z)).
\]
The whole property is not entailed by those
part data. Adding \(R(p,q)\) makes \(W(w)\)
derivable. With the enlarged data but without
the displayed bridge principle, \(W(w)\)
again need not follow.

Thus component properties, relations among
components and a law connecting them to
the whole can play different roles. A label
of emergence does not identify which of
these pieces is missing.
\end{hbtheorem}

The source also emphasizes that an additive
composition rule is itself a principle.
Knowing quantities for several parts does
not, by logic alone, say that the corresponding
quantity for the whole is their sum.
Choosing a different part relation changes
what is being described and need not produce
a monotone enlargement of the same data.
These points concern the relative deductive
test; they do not settle every later
ontological use of the word emergence.

\section{Worked explanatory models and their limits}
\label{sec:p16-dn-examples}

\begin{hbdefinition}{def:p16-dn-thermal-model}{A stipulated two-time-scale expansion model}
To make the source's thermometer illustration
mathematically explicit, suppose two normalized
temperature responses satisfy
\[
 \dot u=a(1-u),\quad \dot v=b(1-v),\quad
 u(0)=v(0)=0,\qquad H(t)=Au(t)-Bv(t),
\]
for \(t\ge0\), with
\[
                     A>B>0,\qquad b>a>0,\qquad Bb>Aa .
\]
The fast response \(v\) enlarges the container;
the slower response \(u\) supplies the eventual
larger expansion. \(H\) is the modeled level
displacement. The equations and parameter
inequalities are assumptions of this
illustration, not empirical constants
asserted for every thermometer.
This real-valued model is outside the
paper's later restricted relational language.
\end{hbdefinition}

\begin{hbtheorem}{thm:p16-dn-thermal-model}{proved=annex:p16.dn-thermal-model}{Initial decrease followed by increase}
The unique differentiable solution gives
\[
 H(t)=A(1-e^{-at})-B(1-e^{-bt}).
\]
It decreases until
\[
                  t_*=\frac{\log(Bb/(Aa))}{b-a}>0,
\]
increases thereafter, and tends to \(A-B>0\).
It has exactly one zero after \(t=0\).
For \(A=2,B=1,a=1,b=3\), that later zero is
\[
                  t_0=\log\frac{1+\sqrt5}{2}.
\]
The derivation shows precisely which response
and composition assumptions produce the
qualitative drop-and-rise pattern.
\end{hbtheorem}

\begin{hbtheorem}{thm:p16-dn-function-and-motive}{proved=annex:p16.dn-function-and-motive}{A benefit or a goal does not supply the missing implication}
The letters in each example stand for distinct atomic claims.
\begin{enumerate}
\item A functional claim \(B\to V\), saying
      that a characteristic contributes a
      viability condition, entails neither
      \(B\) nor \(V\). Even
      \(V\land(V\to(B\lor B'))\) need not
      entail \(B\); an alternative \(B'\)
      can satisfy it.
\item If a current desire \(D\), belief \(K\)
      and opportunity \(O\) are linked to
      action \(A\) by the stipulated rule
      \(D\land K\land O\to A\), then the
      rule and those present conditions
      entail \(A\). They need not entail
      the future achievement \(G\).
\end{enumerate}
Thus a functional consequence is not an
explanation of a characteristic's presence
without further premises, and an explanation
by present motives need not invoke a future
event as a cause. The second rule is an
illustrative conditional, not a universal
empirical law of human behavior.
\end{hbtheorem}

\begin{hbtheorem}{thm:p16-dn-intervention-boundary}{proved=annex:p16.dn-intervention-boundary}{An observational derivation does not select a causal direction}
Let \(U\in\{0,1\}\) have equal probabilities.
Consider the two finite structural models
\[
 \mathcal M_1:\ X=U,\ Y=X,\qquad
 \mathcal M_2:\ Y=U,\ X=Y .
\]
Both have the same observational law, supported
on \((0,0),(1,1)\), and both obey \(X=Y\)
in their ordinary contexts. That relation
permits either variable's ordinary value
to be deduced from the other's.

Replace the equation for \(X\) by \(X=0\)
and retain every other equation.
Then \(\mathcal M_1\) has \(Y=0\) surely,
whereas \(\mathcal M_2\) retains
\(\mathbb P(Y=1)=1/2\).
Consequently the common observational
derivations do not identify the response
to this intervention.
\end{hbtheorem}

The intervention calculation uses the explicit
equation-replacement interpretation of
\cref{def:p16-finite-scm}. It does not treat
the 1948 paper's broad regularity-based use
of causal explanation as a proof of causal
identification.

\section{Using the deductive model as a reference method}
\label{sec:p16-dn-method}

State the exact result sentence and its scope.
Separate the general principles from the
particular conditions. Give the derivation,
identify the principles it needs, and supply
the basic verification witness when applying
the final 1948 model. Record separately why
the intended meanings are appropriate and
what evidence supports the premises.

The model's truth requirement cannot be
replaced silently by a high confirmation
score. Its verification witness is a possible
route, not a recovered observation history.
A successful derivation of one law from a
stronger conjunction does not by itself
establish explanatory unification: the
source explicitly leaves the problem of
distinguishing informative law reduction
from merely adjoining the result unresolved.
These are limits of the identified model,
not permissions to omit premises from
a claimed application.

\subsection{Related reasoning tasks}

\begin{center}
\small
\begin{tabular}{p{.18\linewidth}p{.28\linewidth}p{.42\linewidth}}
\textbf{Task}&\textbf{Inputs}&\textbf{What its result establishes}\\\hline
Logical abduction&
Background, observations and permitted hypotheses&
Added hypotheses are consistent with the background and
entail the observations; see \cref{def:p16-abduction-models}.\\[4pt]
Consistency diagnosis&
System, observations and component normality assumptions&
Minimal withdrawals restore consistency; see
\cref{thm:p16-diag-semantics}.\\[4pt]
1948 deductive explanation&
Generalized \(T\), singular \(C,E\), and an intended interpretation&
A deduction with the stated basic verification witness;
truth-qualified explanation additionally requires the
source's theory and truth conditions.\\[4pt]
Analogical generation&
Source, target and a proposed correspondence&
A candidate mapping or proposal, whose truth needs
separate premises; see \cref{thm:p19-analogy-candidate-truth}.\\[4pt]
Systematic-power scoring&
Sentence \(T\), data \(S\) and a chosen range measure&
A defined conditional ratio of content, without an
automatic truth or adequacy guarantee; see
\cref{thm:p15-power-probability-dual,thm:p15-power-strength-boundary}.
\end{tabular}
\end{center}
